\input BaTeX.tex

\PapierGroesse(66pt,42pt)

\modul{ean,pstex}
\parskip=0pt
\begin{document}
\psSetze{linienstil=}
\begin{psfigur}[scale=2](32,21)
\psPin(-8,-23){
\psCode{0.1 setflat}
\psPolygon[fuellfarbe=0,eckradius=2,linienstil=](8,24)(40,24)(40,44)(8,44)
\psPin(35,40){\psKreis[fuellfarbe=1:0:0,linienstil=](-1,0){2}\psKreis[fuellfarbe=1:1:0,linienstil=](1,0){2}\begin{psGebiet}[fuellfarbe=1:0.5:0,weg=geschlossen]\psKreisbogen(-1,0){2}(-63.43495,63.43495)\psKreisbogen(1,0){2}(116.56505,243.43495)\end{psGebiet}}
\psPin(1,-2){
\psPolygon[fuellfarbe=0.77:0.77:0.75,eckradius=1,linienstil=](10,34)(16,34)(16,38)(10,38)
\psSetze{linienbreite=0.1,linienstil=-}
\psLinie(10,35.333)(12,35.333)
\psLinie(10,36.667)(12,36.667)
\psLinie(14,35.333)(16,35.333)
\psLinie(14,36.667)(16,36.667)
\psLinie(12.2,34.667)(13.8,34.667)
\psLinie(12.2,37.333)(13.8,37.333)
\psKurve(12,34)(12.2,34.667)(12,35.333)
\psKurve(12,36.667)(12.2,36)(12,35.333)
\psKurve(12,36.667)(12.2,37.333)(12,38)
\psKurve(14,34)(13.8,34.66)(14,35.333)
\psKurve(14,36.667)(13.8,36)(14,35.333)
\psKurve(14,36.667)(13.8,37.333)(14,38)
}
\psPin(10,38){
\psPolygon[farbe=1,linienbreite=0.2,linienstil=-](0,0)(4,0)(4,4)(0,4)
\psPolygon[fuellfarbe=1](1.8,0.3)(3.7,0.3)(2.8,2)(3.7,3.7)(1.8,3.7)(0.9,2)
}
}
\psNeueSchicht
\psPin(-8,-22){
\rpin(24,28){\skal[0.48]{\frb{grau7}{\tt 5500 1818 0006 2731}}}
\tpin[r=0](15,41){\skal[0.35]{\frb{weiss}{\sf Schaffhauser}}}
\tpin[r=0](15,38.8){\skal[0.44]{\frb{weiss}{\sf Kantonalbank}}}
}
\end{psfigur}


\end{document}